%!TEX root = ../hutbthesis_main.tex

\chapter{研究理论基础}\label{ux7814ux7a76ux7406ux8bbaux57faux7840}

\section{肌肉骨骼系统建模}\label{ux808cux8089ux9aa8ux9abcux7cfbux7edfux5efaux6a21}

\subsection{人体肌肉骨骼结构与动力学特性}\label{ux4ebaux4f53ux808cux8089ux9aa8ux9abcux7ed3ux6784ux4e0eux52a8ux529bux5b66ux7279ux6027}

人体运动这件事，说到底，就是骨骼、关节和骨骼肌三者在神经系统的统一指挥下，彼此配合完成的一场复杂合作。骨骼像是坚固的杠杆，支撑着整个身体，也给我们一个结构的基础。关节则像运动的枢纽，限制着转动的中心和活动的范围，让每一个动作都处于安全、可控的状态。骨骼肌才是真正的动力来源。神经系统发出信号，肌肉接收到命令后收缩，产生力矩，推动我们的四肢改变姿势、在空间中自由移动。

下肢的肌肉骨骼系统，其实就是我们能站立、走路、跨步、转身时的“主力军”。它包括髋关节、膝关节、踝关节这些关键部位，还有像臀大肌、股四头肌、腘绳肌、胫前肌、腓肠肌在内的几十个肌群。下肢承受的载荷很大，每次落地时地面反冲力也很强，而且要时刻维持身体平衡，控制相当复杂。因此，下肢动力学模型不仅是非线性的，还高度冗余、耦合严重。这里的难点和挑战，生物力学、康复工程、机器人学三大领域都在关注。张天明在研究人体运动生物力学建模时，详细说明骨杠杆、关节约束、肌肉收缩这些环节之间的耦合活性机制，为后续建模打下了坚实的理论基础。

传统的运动控制通常直接用关节角度或关节力矩来控制，虽然这种方法算起来轻松，实施也不复杂，但它根本没考虑到肌肉的驱动特性和生理上的限制。结果就是，动作看起来总是呆板，缺乏自然的流畅感，和真实的人体运动差距很大。肌肉骨骼模型则用肌肉激活作为控制输入，这种方式更贴近人体的神经、肌肉、骨骼之间的控制关系，生成的动作自然也更符合生物运动的规律。

\subsection{神经网络算法}\label{ux795eux7ecfux7f51ux7edcux7b97ux6cd5}

实现"生物合理"运动控制的主要依据是骨肌的力学行为具有严格的生理规律，经典肌力学理论将肌肉运动特性归纳为三大关系，共同决定了肌肉的贡献能力和动态性能。

肌肉长度——紧张感，肌肉只能在一定长度范围内产生最大的张力，过短会使紧张程度急剧下降，肌体的紧张度只能在一定程度。

是肌速——张力关系，肌肉收缩速度与可输出的张力成负相关关系,收缩幅度越紧张越小，肌缩速度越快，筋骨萎缩。

肌肉兴奋物理学是第三种形态，肌肉激活的过程慢，懈怠的过程也有惯性，不能立刻突变，从接受神经信号到产生张力，都明显存在滞后现象。

在仿真环境中，肌肉骨骼控制系统的特性直接通过生理约束引入：关节角度有明确的硬性限制，肌肉激活速率、最大肌力、力矩同样都设有限制。苗展豪\cite{ref27}对上肢骨骼肌肉的力学特性做了细致分析，系统梳理了核心肌群的激活时序以及力的传递规律，为生理约束的量化研究提供了重要参考。只有在这些约束下生成的运动，才能称得上是生物学上合理的运动。这也是本研究与传统虚拟人运动生成方法最大的区别。

\subsection{SMPL 参数化人体模型}\label{smpl-ux53c2ux6570ux5316ux4ebaux4f53ux6a21ux578b}

SMPL是目前计算机视觉、动画、运动捕捉领域里最主流的参数化人体表示模型，它可以使用很少的一组低维参数来准确地表示出不同的体型和姿态的三维人体结构，给运动模仿提供了一个统一、标准的姿态表示。

SMPL 模型包含三类核心参数：

形状参数$\beta$表示个体间的体种差异，是形态参数，如高度、重量、肢体长度、肩宽等，少量的维度就足以涵盖大部分的人体形态差异。

姿态参数$\theta$用轴角来表示，可以完全描述全身关节的朝向和姿态，不会出现欧拉角的万向锁问题。

全局变换参数有根节点（骨盆）三维平移、旋转来确定人体在空间中的位置以及整体的朝向。

SMPL模型在Kinesis系统中扮演运动学奥秘的角色，把运动捕捉数据的解析作为标准关节姿态，获得参考轨迹；将模拟模型输出的关节状态转化为统一的3D关节位置，用于计算误差，实现姿态对齐，并进行量化考核。

\section{运动模仿学习理论}\label{ux8fd0ux52a8ux6a21ux4effux5b66ux4e60ux7406ux8bba}

（1）运动数据的表示与时序特征结构

运动模仿的精髓在于对一段连续时序的运动进行学习和再现，模仿学习的前提条件是高质量的运动数据，其表现方式影响了学习效果，也影响了考核的准确度。完整的运动数据一般是由多个层次的信息组成的：

底层关注关节的旋转信息，用欧拉角、四元数或轴角描述，这些数据直接反映关节本身怎么转动。中间层包含关节在三维空间中的位置信息，展示肢体运动的轨迹。再往上，是速度、角速度和加速度，体现运动的动态特征和流畅度。最高层是全局状态，也就是根节点的情况，包括骨盆的位置、速度、高度和朝向，这些因素决定了整个运动的轨迹和身体的平衡稳定。

这些信息按时间顺序依次排列，形成时序序列，具有高维强连续性，运动模仿学习空间姿态外，更重要的是学习时间演进规律，因此属于高维度时令决策问题的典型。

（2）运动模仿学习的主流范式与特点

运动模仿学习是使智能体再现示范动作的外在特征、节奏、力度和动态等，根据学习目标的不同可以分为三类范式：

第一是学习照搬照抄，以轨迹为依据，把复现示范轨迹作为主要目的，将预测姿态位置错误直接最小化，追求"看起来差不多"。

第二是动力模仿学习，对肌肉活动性、关节力矩、地面反作用力等恐龙特性的研究，既考虑了表象，又考虑运动的内在机理。

第三是意图模仿学习，不会精雕细琢每一帧，而是懂得达到目标点、跨越关口、保持平衡等示范动作所要达到的目的。

本文把轨迹模仿和生物力学约束结合起来，保证运动外观逼真的同时又保证肌肉激活、关节运动符合生理规律，达到形似而合理的控制目的。

（3）运动数据预处理与标准化流程

如果直接使用这些数据进行训练和评价，会使指标波动比较大，结论不一致，实验难以重复，在原始运动抓手数据中会出现噪音、漂移、克隆率不一、高度偏移，朝向杂乱、异常跳变等问题。

本研究采用的预处理流程包括：

噪声滤波用高斯滤波和滑动平均滤波直接干掉高频抖动和异常值，运动看起来就平滑不少。帧率统一其实就是把来自不同来源的数据，不同帧率的都拉到一个标准上，重采样，确保时序对齐，大家都在同一个节奏里。根节点高度校正很重要，自动调整垂直位移，让足部和地面紧密接触，别出现悬空或者穿模的情况。朝向归一化也不能少，把所有运动的初始方向统一起来，这样朝向差异不会影响评价结果。时间序列对齐和裁剪处理，解决了运动序列长度不同的问题，让它们都能用同一个标准去度量。数据增强的部分，用随机旋转、加点小噪声、提高速度等方法来强化策略的适应性。标准化之后，运动数据统一到同一个尺度、空间和时间基准，实验的稳定性、公平性和可重复性都提升了不少。

\section{强化学习与深度强化学习理论}\label{ux5f3aux5316ux5b66ux4e60ux4e0eux6df1ux5ea6ux5f3aux5316ux5b66ux4e60ux7406ux8bba}

（1）强化学习的基本框架与马尔可夫决策过程

加强学习，是依靠与环境互动、试错、获取反馈、不断完善策略，完成诸如肌肉、骨体控制等复杂动态决策问题的机器学习。加强学习可以建模为马尔可夫决策流程（MDP），主要因素有5个，分别是状态、动作、状态转移、奖励和折扣因子。

状态空间$s$包含系统所有的可观测信息，即关节角度、速度、肌肉激活、参考轨迹、根节点位置等。

动作空间$a$是智能体所发出的控制指令，也就是肌肉激活信号。

状态转移描述当前动作执行后，系统转移到下一状态的概率分布，体现环境的动态特性。

奖励函数$r$是对于当前动作好坏的一种量化反馈，从而促使策略朝着高精度、高稳定、生物合理方向发展。

折扣因子$\gamma$用于权衡短期奖励与长期奖励，收益在0-1之间，保证回报的稳定性。

Actor（策略网络）$\pi$作为映射函数，承担着从状态到动作的生成任务，其核心是运用神经网络对控制动作进行精准拟合。

Critic（价值网络）$V(s)/Q(s,a)$是当前状态或者状态-动作对的长期收益的评价，用以稳定策略更新，降低训练波动。

Actor不断地和仿真环境交互来减小跟踪误差、保证平衡以及满足约束，从而得到高质量的运动控制策略。

（2）深度强化学习架构与优势

深度强化学习（DRL）把深度学习的高维度特征提取能力与加强学习决策能力结合起来，能够处理这种高维、非线性、强耦合、连续动作空间筋骨系统复杂的问题。主流架构有三类：

基于价值的方法以DQN为代表，用价值网络间接得到策略，适合于离散动作空间。

于连续控制之情境颇为适配、直接针对策略函数施以梯度上升优化的策略梯度法，却存在训练不稳定的弊端。

Actor-Critic架构将两者优势结合起来，利用Actor网络产生动作，运用Critic网络评价动作的价值，提高训练稳定性和衔接速度。

（3）近端策略优化算法（PPO）原理与优势

PPO是目前机器人控制、肌肉骨骼仿真、运动模仿中应用最广泛、最稳定的策略梯度算法，也是
Kinesis 系统的核心训练算法。

PPO的核心创新是在一个更新周期内，防止战略变化过大，导致训练失败，动作不稳定，模式发散，而使用裁剪替代目标函数的限制策略更新幅度，PPO使用多次经验数据，以提高样本量降低训练成本。

PPO三个突出的优点,训练很稳定，适合肌肉骨料这样复杂的动力系统，超参数不敏感，化繁为简，不需要复杂调试，实验标准化规范化。

本文用轨迹模仿误差作为主要的优化目标，加上生物合理性惩罚项，最后得到一种既能准确模仿，又能满足生理约束的肌肉激活策略。

PPO算法具有显著的训练稳定性和样本利用率优势，在这种情况下策略更新保守、探索不足、梯度波动大等问题仍然存在于传统版本的高维连续动作空间、肌肉骨头等复杂运动系统。

先说策略更新。我们把裁剪阈值做成自适应的，让它能根据训练进度自动调整。训练刚开始的时候，阈值放宽一点，方便探索；等到后期再慢慢收紧，保证策略稳定，不会出现更新特别慢或者一下子突变的情况。

再来看优势估计。这里用上了加权广义优势估计法。它对多步时序误差进行加权，权重取决于每一步的可信度。这样做能有效减少高方差样本对结果的干扰，让优势估计整体更平滑、更可靠。

最后，策略网络结构也做了优化。输出层加入了状态依赖的方差，让策略探索的强度能跟着状态的不确定性自动调整。这样就不会无效地瞎探索，训练效率明显提升。

这些改进让算法在肌肉骨骼高维、强耦合、易崩溃的复杂环境中，兼顾探索能力、稳定性、收敛速度，为后续运动模仿任务提供更稳健的优化基础。

\section{运动模仿量化评估体系}\label{ux8fd0ux52a8ux6a21ux4effux91cfux5316ux8bc4ux4f30ux4f53ux7cfb}

（1）平均每关节位置误差（MPJPE）

MPJPE运动仿真界最通用的测关节位置量测和欧陆平均距离位置的精度指标，其计算公式为：

\begin{equation}
\mathrm{MPJPE}=\frac{1}{N}\sum_{i=1}^{N}\left\|\widehat{x}_{i}-x_{i}\right\|_{2}
\tag{2-1}
\end{equation}

其中$N$是考核关节总数，\({\widehat{x}}_{i}\)是预测关节位置的模型，也是真正的参照部位，MPJPE越小，代表愈精准、愈精确，动作"越像"。

（2）Frame coverage

帧覆盖率指模型在稳定运行状态下，即不崩溃、不失控、不异常终止时，所成功执行的帧数占总帧数的比例。

\begin{equation}
\operatorname{Coverage}=\frac{N_{\mathrm{valid}}}{N_{\mathrm{total}}} \times 100\%
\tag{2-2}
\end{equation}

Coverage：帧覆盖率，即运动连续稳定度。\(N_{\mathrm{valid}}\)为成功、没有崩溃、没有跌倒、关节不超限有效帧数。\(N_{\mathrm{total}}\)：整个动作序列总帧数。含义为跑到底的帧数所占的百分比，越高说明运动越稳定，不容易出现崩溃的情况。

（3）任务成功率（Success Rate）

成功率是指能够完全执行动作、保持平衡、不穿模的序列在多次重复测试中所占的比例,是鲁棒性指标中最最贴近工程应用的一种。
\begin{equation}
\operatorname{SuccessRate}=\frac{N_{\mathrm{success}}}{N_{\mathrm{total}}} \times 100\%
\tag{2-3}
\end{equation}

SuccessRate：任务成功率，衡量整体鲁棒性、可用性。\(N_{\mathrm{success}}\)是完整跑完、不倒、不穿模、不超限的序列数\(N_{\mathrm{total}}\)：总共测试的动作序列数量。含义为完成任务次数占总次数的比例，越大表示策略越可靠、工程价值越高。

（4）组合式量化评估体系

单项指标不能对运动模仿质量进行全面客观科学的评估，MPJPE只能反映准确性，不能反映稳定性，成功率只能反映成功或失败，行动质量不能体现出来，帧覆盖率只能体现连续性，而不能反映准确性。

本文提出综合评价体系，即MPJPE+帧覆盖率+成功率，MPJPE回答动作精度的精确程度。帧覆盖率回答动作能稳定到什么程度，成功率回答动作能不能成功。

\section{Kinesis 仿真平台与 MuJoCo 物理引擎}\label{kinesis-ux4effux771fux5e73ux53f0ux4e0e-mujoco-ux7269ux7406ux5f15ux64ce}

（一）Kinesis 平台架构与设计理念

Kinesis
专为下肢肌肉骨骼运动的模仿和生物合理控制而设计，是面向强化学习研究的一个完整框架。它追求生理可信、运动自然、流程标准化和实验结果可复现。这个平台把数据处理、运动学计算、策略学习和标准化评价这些核心功能打包在一起，让肌肉骨骼运动控制算法从数据输入、模型训练、仿真测试到量化评价，全流程都有了有力支撑。在这个环境下，可以更高效地开发、验证和比较不同的控制策略，推动生物力学和运动科学的深入发展。

Kinesis整体采用模块化、低耦合、高扩展架构，主要由四大功能模块组成：

（1）数据处理模块：负责运动抓拍数据的加载、解析、清洗、对齐、滤波、增强等工作，负责对运动捉拍资料的分析、清理、正向、过滤波和增强等。支持多种运动数据格式的统一转换，能够自动完成噪声的去除、时间的同步、空间的归一化和数据的拓展，确保输入数据规范可靠，为后来学习提供了高质的运动序列。

（2）运动学计算模块：可快速得到关节姿势示意、三维关节位置计算、运动轨迹重建、误差分析，利用SMPL参数化人体模型完成高效批量前向运动求解。

（3）战略学习模块：支持肌肉激活策略的端对端训练，以PPO近端策略优化为核心。在灵活的网络结构配置中，在高维度、连续动作空间中稳定地进行肌肉活化模式的学习，产生生物合理的运动策略，根据任务需求自由地改变策略网络、价值网络以及奖励函数的设计。

（4）标准化考核模块：内含标准化评价流程，输出三个主要的量化指标MPJPE、帧覆盖率、任务成功率，可以消除主观评价差异，使不同算法之间横向比较公平合理，评价过程固定，参数公开，结果可重复使用。

Kinesis的显著优势在于流程固定、配置透明、参数可结果可复现，能够严格保证实验可重复，很好地满足了本次研究的核心要求——标准化、可量化肌肉骨骼运动模仿。

（二）MuJoCo 物理引擎与动力学仿真支撑

MuJoCo高精度物理仿真引擎，面向多关节和接触、高复杂度系统，广泛应用于生物力学、机器人学、强化学习等领域，具有高稳定性、高精度、高效率三大优势，是理想的肌肉骨骼运动仿真物理引擎。

MuJoCo专为生物力学可信仿真而设计，其主要功能有如下几个方面：

（1）能计算多刚体、肌肉驱动和接触动力学以及肌肉收缩弧、关节力矩传递和复杂的接触行为，并进行大范围内的模拟仿真，保证对多刚体、肌体驱动和接触动力学等问题的高效解决。

（2）对地面摩擦、碰撞反应、关节阻尼、脂肪约束、重力效应、地面反作用力等进行高精度的物理特征仿真使其尽可能真实，进一步保证仿真动力学和真实的人体生物力学一致。

（3）灵活的生理限制设置能够准确模拟人体肌肉骨系统的生理约束，保证运动的生物合理性，如肌肉执行器参数、关节活动范围、力矩限制限制消耗限制等，均可以自定义。

（4）数值稳定性高，适用于长期、大尺度、多轮次的加强学习培训，用先进的数值积分算法，以低误差、低崩溃率，在同等条件下获得较高精度，保证了强化学习和培训的稳定结果。

MuJoCo赋予肌肉骨模型高稳定性的物理模拟，保证仿真运动接近真实的人体行为，并给予训练、测试、验证生物合理的肌肉骨骼运动控制算法，以坚实的基础。

（三）平台适用性与本研究的契合性

Kinesis平台与这一研究需求相契合，能够全面支持标准化、可反复、量化筋骨运动模仿研究，从模型、流程、评价三个方面进行，主要体现在以下三点：

（1）原生支持肌肉骨骼模型和生物合理性约束，不需要从无到有地建立复杂的肌肉骨骼模型，平台自有生理可信的肌肉骨骼模型，可以对关节角度、肌肉活化、能耗、运动通畅等各个方面进行多层次的生物约束。

（2）平台不会因不同评价标准而造成结果的不可比较性，可以自动生成MPJPE、帧覆盖率和任务成功率等指标，保证本研究评价客观公正、可复制。

（3）固定随机种子、固定测试集、固定超参数，以及完全可以复现的实验环境，保证每次训练、测试结果稳定一致，解决目前领域研究难以复现、结论不可靠的痛点，为本研究提供科学、严谨、可验证的实验基础。
